但生态系统的恢复并非总是朝着良性方向发展。部分水域在禁渔后出现鱼类种群数量增长过快、过度摄食水生植物和浮游动物的现象，进而导致水体浑浊、底栖植被大面积减少，部分水域甚至形成大面积"水下荒漠"。与此同时，长江江豚种群数量已回升至约1249头，2024年长江鲟自然繁殖产卵场重新被发现；而鳄雀鳝、巴西龟等外来入侵物种已在鄱阳湖等流域支流形成稳定种群。十年禁渔政策确实取得了显著的资源恢复成效，但恢复后的生态系统未必更加健康，营养级失衡、江湖连通性恢复不足以及外来物种入侵扩散等问题依然突出。因此有必要构建统一的动力学模型，将"资源恢复"与"系统健康"两个概念进行区分并分别评估。\cite{report2022,report2023}。
